\section{K2 技术决策蓝图：四个目标如何互相牵制}

上一章说明 K2 不是单点优化，本章把这些目标整理成决策蓝图，回答“模型公司到底在同时优化什么”。基础模型、Agentic 能力、开源生态和商业化不是并列口号，而是互相牵制的四个目标。基础模型越强，Agent 产品和 API 越有底气；开源越彻底，生态扩散越快，但商业捕获越难；Agentic 能力越强，工具和反馈越重要，训练和评测也越复杂。

\begin{knowledgebox}{K2 四目标矩阵}
\begin{tabularx}{\textwidth}{p{0.18\textwidth}p{0.30\textwidth}p{0.38\textwidth}}
\toprule
目标 & 要优化什么 & 会牵动什么 \\
\midrule
Base Model & 通用能力、代码能力、推理能力 & 数据质量、token efficiency、优化器和算力。 \\
Agentic & 多轮工具使用、环境反馈、任务泛化 & RL、测试时计算、工具协议和评测。 \\
Open Source & 外部验证、开发者生态、信任 & 商业捕获、支持成本和竞品复用。 \\
Business & API、产品、企业和开发者付费 & 模型边界、Agent 应用边界和成本结构。 \\
\bottomrule
\end{tabularx}
\end{knowledgebox}

\begin{importantbox}{为什么这四个目标不能拆开看}
如果只有强基座而没有 Agentic 能力，模型仍然像“缸中之脑”；如果只有 Agent 产品而基座不强，泛化会受限；如果开源但没有商业路径，生态难以持续；如果商业化过早，又可能限制研究和开源扩散。
\end{importantbox}

\subsection{训练效率与产品能力的循环}

四目标矩阵里最容易被低估的是训练效率，本节说明它为什么会传导到产品。Token efficiency 看起来是训练指标，但它会影响产品能力。更高的 token efficiency 意味着同样数据和算力下得到更强底座；更强底座让 Agent 更能处理复杂任务；复杂任务产生更多反馈和使用数据；这些反馈又可能进入下一轮训练和产品迭代。这个循环如果跑起来，模型公司就不只是训练模型，而是在构建能力飞轮。

\begin{knowledgebox}{能力飞轮}
\begin{tabularx}{\textwidth}{p{0.20\textwidth}p{0.34\textwidth}X}
\toprule
环节 & 作用 & 风险 \\
\midrule
训练效率 & 降低获得能力的成本 & 优化器或数据策略不稳定。 \\
Agent 产品 & 让模型进入真实任务 & 失败率、工具安全和用户信任。 \\
反馈数据 & 记录真实任务中的成功和失败 & 噪声大、隐私和选择偏差。 \\
下一轮训练 & 用反馈修正模型和工具策略 & 评测不准会放大错误。 \\
\bottomrule
\end{tabularx}
\end{knowledgebox}

\subsection{本章小结}

K2 的技术决策不是单点最优，而是四目标联合优化。理解这个蓝图，才能理解为什么访谈会在 token efficiency、Agentic 泛化、开源、API 和商业模式之间来回切换。

